%!TEX root = ./main.tex

\usetikzlibrary{positioning,fit,backgrounds}

\definecolor{ranblue}{RGB}{46,91,186}
\definecolor{coreorange}{RGB}{222,122,36}
\definecolor{ranpurple}{RGB}{109,53,151}
\definecolor{softgreen}{RGB}{73,139,74}
\definecolor{softgray}{RGB}{92,98,105}
\definecolor{attackred}{RGB}{178,44,44}

\newcommand{\source}[2]{%
  \par\vspace{0.25em}%
  {\tiny\textcolor{gray}{\href{#1}{\underline{#2}}}\par}%
}

% Inline expansion of an acronym at first use: NTN \expand{Non-Terrestrial Network}
\newcommand{\expand}[1]{\textcolor{softgray}{(#1)}}

% Small gray glossary line, for slides with several acronyms at once.
\newcommand{\glossline}[1]{%
  \par\vspace{0.4em}%
  {\scriptsize\color{softgray}\raggedright#1\par}%
}

\newcommand{\plainterm}[2]{%
  \textbf{#1}\hspace{0.25em}{\small\textcolor{softgray}{#2}}%
}

\newcommand{\ask}[1]{%
  \begin{tcolorbox}[colback=blue!4,colframe=ranblue,boxrule=0.7pt,arc=2pt,
    left=5pt,right=5pt,top=4pt,bottom=4pt]
    \centering\large #1
  \end{tcolorbox}%
}

% Font is deliberately \normalsize, not \large: at \large a one-sentence takeaway
% wraps to two lines and the box grows into whatever is below it.
\newcommand{\takeaway}[1]{%
  \begin{tcolorbox}[colback=orange!7,colframe=coreorange,boxrule=0.7pt,arc=2pt,
    left=5pt,right=5pt,top=4pt,bottom=4pt]
    \centering\normalsize #1
  \end{tcolorbox}%
}

\tikzset{
  netnode/.style={draw=ranblue,fill=blue!5,rounded corners=3pt,
    minimum width=1.9cm,minimum height=0.82cm,align=center,font=\bfseries},
  corenode/.style={draw=coreorange,fill=orange!7,rounded corners=3pt,
    minimum width=1.9cm,minimum height=0.82cm,align=center,font=\bfseries},
  oranode/.style={draw=ranpurple,fill=purple!6,rounded corners=3pt,
    minimum width=1.9cm,minimum height=0.82cm,align=center,font=\bfseries},
  attacknode/.style={draw=attackred,fill=red!6,rounded corners=3pt,
    minimum width=1.9cm,minimum height=0.82cm,align=center,font=\bfseries},
  flow/.style={-{Latex[length=2.4mm]},thick,draw=softgray}
}

% Full-width question/answer block: spans the whole frame, no box.
\newcommand{\qa}[2]{%
  \par\vspace{0.55em}%
  {\large\raggedright
    \noindent\hangindent=1.25em\hangafter=1
    \textcolor{ranblue}{\textbf{Q:}}\hspace{0.35em}\textbf{#1}\par
    \vspace{0.35em}
    \noindent\hangindent=1.25em\hangafter=1
    \textcolor{coreorange}{\textbf{A:}}\hspace{0.35em}#2\par}%
}

% ---------------------------------------------------------------------------
% Macros specific to the Class 5 "seven papers" lecture.
% ---------------------------------------------------------------------------

% Citation card for the paper a segment is built on. Use ONCE per paper, on the
% frame that introduces it. #1 title, #2 authors/affiliation (short), #3 venue,
% #4 URL (DOI or official page).
\newcommand{\papercard}[4]{%
  \begin{tcolorbox}[colback=purple!4,colframe=ranpurple,boxrule=0.7pt,arc=2pt,
    left=6pt,right=6pt,top=4pt,bottom=4pt]
    \raggedright
    {\small\textbf{#1}\par}
    \vspace{0.15em}
    {\footnotesize #2\par}
    \vspace{0.15em}
    {\footnotesize\textcolor{softgray}{#3}\quad
      {\scriptsize\href{#4}{\textcolor{ranblue}{\underline{paper}}}}\par}
  \end{tcolorbox}%
}

% Link to the official talk video, with a suggested clip length for class.
% #1 URL, #2 label (e.g., "SIGCOMM'25 talk, play 0:00-1:30").
\newcommand{\watch}[2]{%
  \par\vspace{0.25em}%
  {\small\textcolor{ranblue}{\raisebox{-0.1em}{\(\blacktriangleright\)}\;\href{#1}{\underline{#2}}}\par}%
}

% The recurring "assumption -> broken by" pair that stitches the lecture
% together. #1 the hidden assumption, #2 who breaks it.
\newcommand{\brokenassumption}[2]{%
  \par\vspace{0.4em}%
  {\raggedright
    \noindent\textcolor{softgray}{\textbf{Assumption:}}\hspace{0.35em}#1\par
    \vspace{0.25em}
    \noindent\textcolor{attackred}{\textbf{Broken by:}}\hspace{0.35em}\textbf{#2}\par}%
}
